\documentclass[11pt]{article}
\usepackage{amsmath,amssymb,amsthm}
\usepackage[margin=1.1in]{geometry}
\usepackage{hyperref}

\newtheorem{proposition}{Proposition}
\newtheorem{remark}{Remark}
\newcommand{\xbar}{\bar{x}}
\newcommand{\phat}{\hat{p}}

\title{Bernoulli Maximum-Likelihood Estimation, Part I:\\
The Theoretical Solution and the Estimator's Report Card}
\author{Yuval Lavie}
\date{\today}

\begin{document}
\maketitle

\begin{abstract}
Pure statistics --- no losses, no descent. From i.i.d.\ samples of a
Bernoulli distribution we derive the maximum-likelihood estimator in
closed form, then judge it the way a statistician judges any estimator:
bias, variance, efficiency against the Cram\'er--Rao bound, mean squared
error, consistency, asymptotic normality (hence confidence intervals),
and sufficiency. The companion script \texttt{mle\_theoretical.py}
verifies the closed form against a brute-force search and every
quantitative property by Monte Carlo over thousands of replicated
experiments. The translation of this problem into a machine-learning
loss, and its solution by gradient methods, is Part~II
(\texttt{mle\_practical.tex}).
\end{abstract}

\section{From the likelihood to the estimator}

Let $x_1,\dots,x_n \overset{\text{iid}}{\sim} \mathrm{Bernoulli}(p)$ with
$x_i \in \{0,1\}$ and unknown $p \in (0,1)$. The question of maximum
likelihood: which $p$ makes the observed sample most probable?

\paragraph{Step 1: the likelihood is a product of two kinds of factors.}
By independence, the probability of observing exactly this sample is
\begin{equation}
  L(p) \;=\; \prod_{i=1}^{n} \Pr(x_i)
       \;=\; \prod_{i=1}^{n} p^{x_i}(1-p)^{1-x_i}.
\end{equation}
Each factor is just a case split written as one expression: it equals $p$
when $x_i = 1$ and $1-p$ when $x_i = 0$. So if the sample contains
$m = \sum_i x_i$ ones and $n - m$ zeros, collecting the factors gives
\begin{equation}
  L(p) \;=\; p^{\,m}\,(1-p)^{\,n-m}.
  \label{eq:likelihood}
\end{equation}
Already here the data enters only through the count $m$: how the ones and
zeros are ordered is irrelevant. (This observation becomes
\emph{sufficiency} in Section~\ref{sec:reportcard}.)

\paragraph{Step 2: the logarithm turns the product into a sum.}
$\log$ is strictly increasing, so it preserves the location of a maximum:
$\arg\max_p L(p) = \arg\max_p \log L(p)$. Applying it to
\eqref{eq:likelihood},
\begin{equation}
  \ell(p) \;=\; \log L(p) \;=\; m\log p + (n-m)\log(1-p).
  \label{eq:loglik}
\end{equation}
This is the whole reason for working in logs: powers become products,
products become sums, and sums are easy to differentiate (and do not
underflow numerically the way a product of $5000$ numbers in $(0,1)$
would).

\paragraph{The closed form.}

\begin{proposition}
Assume $0 < m < n$. Then $\phat_{\mathrm{MLE}} = m/n = \xbar$, and it is
the unique maximizer of $\ell$ on $(0,1)$.
\end{proposition}

\begin{proof}
Differentiate \eqref{eq:loglik} term by term:
\begin{equation*}
  \ell'(p) \;=\; \frac{m}{p} - \frac{n-m}{1-p}.
\end{equation*}
Set it to zero and clear denominators:
\begin{equation*}
  \frac{m}{p} = \frac{n-m}{1-p}
  \;\;\Longrightarrow\;\;
  m(1-p) = (n-m)\,p
  \;\;\Longrightarrow\;\;
  m - mp = np - mp
  \;\;\Longrightarrow\;\;
  m = np,
\end{equation*}
hence
\begin{equation}
  \boxed{\;\phat_{\mathrm{MLE}} \;=\; \frac{m}{n} \;=\; \xbar\;}
  \label{eq:phat}
\end{equation}
--- the fraction of ones in the sample. Uniqueness:
$\ell''(p) = -m/p^{2} - (n-m)/(1-p)^{2} < 0$, so $\ell$ is strictly concave
on $(0,1)$ and the stationary point is the global maximum.
\end{proof}

\begin{remark}[Edge cases]
If $m = 0$ or $m = n$, every factor of the likelihood pulls in the same
direction: $L$ is monotone on $(0,1)$ and its supremum sits at the
boundary, which is again $m/n$ --- the formula survives.
\end{remark}

\section{The estimator's report card}
\label{sec:reportcard}

An estimator is a \emph{random variable}: a new sample produces a new
$\phat$. So $\phat$ has a distribution of its own, and we judge it by
that distribution's shape. Throughout, use $m \sim \mathrm{Binomial}(n,p)$,
so $\mathbb{E}[m] = np$ and $\mathrm{Var}(m) = np(1-p)$.

\paragraph{Bias.} $\mathbb{E}[\phat] = \mathbb{E}[m]/n = p$: the
estimator is \emph{unbiased} --- over repeated samples it is centered on
the truth exactly, at every $n$.

\paragraph{Variance and standard error.}
\begin{equation}
  \mathrm{Var}(\phat) \;=\; \frac{\mathrm{Var}(m)}{n^{2}}
  \;=\; \frac{p(1-p)}{n},
  \qquad
  \mathrm{SE} \;=\; \sqrt{\frac{p(1-p)}{n}},
  \label{eq:var}
\end{equation}
estimated in practice by plugging $\phat$ in for $p$. The $1/n$ decay is
the rate at which data buys precision.

\paragraph{Efficiency: the Cram\'er--Rao bound is attained.}
The Fisher information of a single Bernoulli observation is
\begin{equation*}
  I(p) \;=\; \mathbb{E}\!\left[-\frac{d^{2}}{dp^{2}}
  \log f(X;p)\right]
  \;=\; \mathbb{E}\!\left[\frac{X}{p^{2}} + \frac{1-X}{(1-p)^{2}}\right]
  \;=\; \frac{1}{p} + \frac{1}{1-p}
  \;=\; \frac{1}{p(1-p)},
\end{equation*}
and the Cram\'er--Rao inequality states that \emph{every} unbiased
estimator $T$ of $p$ obeys $\mathrm{Var}(T) \ge 1/\bigl(n I(p)\bigr)$
(derivation: \texttt{Theory/statistics}). Here
\begin{equation}
  \frac{1}{n\,I(p)} \;=\; \frac{p(1-p)}{n} \;=\; \mathrm{Var}(\phat):
  \label{eq:crlb}
\end{equation}
the bound holds with \emph{equality}. The MLE is \emph{efficient} ---
provably unimprovable among unbiased estimators, at every finite $n$.

\paragraph{Mean squared error.} The general decomposition
$\mathrm{MSE} = \mathrm{Var} + \mathrm{Bias}^{2}$ collapses to
$\mathrm{MSE}(\phat) = p(1-p)/n$: for unbiased estimators, accuracy is
purely a variance story.

\paragraph{Consistency.} Unbiasedness plus vanishing variance forces
$\phat \to p$ in probability (Chebyshev). Since $\phat = \xbar$, this
\emph{is} the law of large numbers, specialized to coin flips.

\paragraph{Asymptotic normality, and what it buys.}
By the central limit theorem,
\begin{equation}
  \sqrt{n}\,(\phat - p) \;\xrightarrow{\;d\;}\;
  \mathcal{N}\bigl(0,\; p(1-p)\bigr),
  \label{eq:clt}
\end{equation}
so for large $n$ the sampling distribution of $\phat$ is approximately
$\mathcal{N}(p,\, p(1-p)/n)$. This is what buys \emph{confidence
intervals}:
\begin{equation}
  \phat \;\pm\; 1.96\,\sqrt{\frac{\phat(1-\phat)}{n}}
  \label{eq:ci}
\end{equation}
covers the true $p$ in about $95\%$ of repeated experiments. (The
asymptotic variance $1/I(p)$ appearing here is no accident --- it is the
general asymptotic behavior of MLEs.)

\paragraph{Sufficiency.} The likelihood \eqref{eq:likelihood} depends on
the data only through $m$ --- the factorization criterion --- so $m$ is a
\emph{sufficient statistic}: compressing $n$ observations into one count
loses nothing about $p$.

\begin{remark}[Invariance, used in Part~II]
The MLE of any function $g(p)$ is $g(\phat)$. This is why Part~II may
reparameterize the problem through a sigmoid and still recover the same
estimator.
\end{remark}

\section{Experiment}

With $p = 0.3$ and $n = 5000$ (seed 7): $m = 1489$, so
$\phat = 0.297800$, agreeing with a brute-force maximization of
\eqref{eq:loglik} over a $10^{5}$-point grid. The report card, verified
by Monte Carlo over $R = 5000$ replicated experiments: empirical mean of
$\phat$ is $0.299910$ against $p = 0.3$ (unbiasedness); empirical
variance $4.244\times10^{-5}$ against
$p(1-p)/n = 4.200\times10^{-5}$, which equals the Cram\'er--Rao bound
(efficiency); the $95\%$ interval \eqref{eq:ci} is $[0.2851, 0.3105]$
and contains the truth. The figure shows, left to right: the true and
fitted probability mass functions side by side; the log-likelihood
\eqref{eq:loglik} with its maximum at \eqref{eq:phat}; and the sampling
distribution of $\phat$ --- the histogram of the $5000$ replicated
estimates --- with the CLT normal \eqref{eq:clt} drawn over it,
indistinguishable.

One sentence to keep: the estimator is a random variable, and everything
we proved --- unbiased, efficient, normal --- is a statement about
\emph{its} distribution, not about any single number.

\end{document}
